\Th (Великая теорема Ферма при $n = 3$)
	Уравнение $x^3 + y^3 = z^3$ не имеет нетривиальных целочисленных решений.
\\
\\
В следующих трех билетах мы будем доказывать теорему Ферма. 
\\
Сразу сделаем несколько замечаний. Прежде всего, удобнее
говорить о нетривиальных решениях в целых числах (то есть о таких, в которых ни одно число не равно нулю). Тогда можно
говорить как об уравнении $x^3 + y^3 = z^3$, так и об уравнении $x^3 + y^3 + z^3 = 0$ Из последнего следует, что переменные можно менять местами друг с другом. \\
Давайте сначала найдём идею доказательства, а потом его строго запишем.
\\
Мы хотим использовать числа Эйзенштейна для доказательство этого факта. Почему? Дело в том, что в числах Эйзенштейна есть кубический корень из 1. Поэтому мы можем разложить $x^3 + y^3$ на линейные сомножители. А именно, $x^3 + y^3 = (x + y)(x + \omega y)(x + \omega^2y)$  (это равенство можно проверить непосредственно). \\
Итак, мы будем доказывать, что уравнение
$x^3 + y^3 = z^3$ не
имеет нетривиальных решений в числах Эйзенштейна. Как и в пифагоровых тройках, можно предполагать попарную взаимную
простоту x, y, z. Мы имеем равенство:
\begin{center}
    $(x + y)(x + \omega y)(x + \omega^2y) = z^3$
\end{center}
Если продолжать аналогию с пифагоровыми тройками, надо выяснить взаимную простоту трёх сомножителей. Заметим, что
$(x + y, x + \omega y) = ((1 - \omega)y, x + y)$.
Что за число $1 − \omega$? Его норма, как легко посчитать, равна 3,
значит, это простое число Эйзенштейна. Для сокращения записи обозначим
$\lambda = 1 - \omega$. Вернёмся к наибольшему общему делителю $x + y$ и $x + \omega y $.
\\
\\
\Lemma
 Если $\lambda$ делит $x^3 + y^3$ для взаимно-простых $x, y$, то $(x + y, x + \omega y) = (x + y, x + \omega^2 y) = (x + \omega^2y, x + \omega y) = \lambda$
 \\
 \Proof
Так как y и x + y не имеют общих делителей, то
$(x + y, x + \omega y) = (\lambda y, x + y) = (\lambda, x + y) = \lambda$ или 1.
\\
Далее $(x + y, x + \omega^2 y) = (x + y,(1 - \omega^2)y) = (x + y, \lambda y) = (x + y, x + \omega y)$,
\\
$(x + \omega y, x + \omega^2y) = (x + \omega y, \lambda \omega y) = (x + \omega y, \lambda y) = (x + y, x + \omega y)$.
То есть либо все три сомножителя попарно взаимно просты, либо у каждой пары наибольший общий делитель равен $\lambda$. Первый случай невозможен, так как по условию $\lambda$ делит $x^3 + y^3$
\EndProof
\\
\\
\Lemma
	$z \in \Z[\omega]$ либо делится на
    $\lambda$, либо сравним с обратимым элементом по модулю 3. Более того, если $\lambda \nmid z$. Тогда $z^3 \equiv_9 \pm 1$.

\Proof
Пусть x = $a+b\omega$. По модулю 3 можно представить x как $a_0+b_0\omega$, где $a_0, b_0$ $\in$ $\{0, −1, 1\}$. Из 9 вариантов 6 задают
обратимый элемент, а оставшиеся 3 - это $0, 1 -\omega = \lambda, -1 + \omega$ =
$-\lambda$, то есть делятся на $\lambda$. 
\\
\\
С другой стороны, пусть $x = 3w + u$, где u — обратимый элемент. Тогда
$(3w + u)^3 = 27w^3 + 27w^2u + 9wu^2 + u^3  \equiv_3 u^3  \equiv_9 \pm1$.
\EndProof
\\
\\
\Lemma
	Если $\pm x^3 \pm y^3 \pm z^3 = 0$ для некоторых $x, y, z \in \Z[\omega]$, то $\lambda \mid xyz$.
\Proof
	Если это не так, то $0 = \pm x^3 \pm y^3 \pm z^3 \equiv_9 \pm 1$ --- противоречие.
\EndProof
